
\subsubsection{2.1 AI-to-Human Alignment: Well-Measured, One-Directional}

The dominant paradigm in AI alignment research operationalizes ''alignment'' as the degree to which AI outputs match human preferences, values, or intent. RLHF (Ouyang et al., 2022; Bai et al., 2022) provides preference-labeled training data from human raters. Benchmarks such as HELM \citep{Liang2022Holistic}, TruthfulQA \citep{Lin2022TruthfulQA}, and BIG-Bench \citep{Srivastava2022Beyond} measure AI performance on tasks that proxy human-valued behaviors. This literature is mature, with standardized evaluation infrastructure and reproducible scores across model generations.

A critical feature of this paradigm is that it treats human evaluators as a fixed, stable signal source. The implicit assumption is that rater behavior is stationary: raters in 2024 apply the same judgment standards as raters in 2022. Whether this assumption holds as raters accumulate experience with AI-generated text has not been systematically examined.

\subsubsection{2.2 Human Behavioral Adaptation: Qualitative Evidence, No Measurement Framework}

A smaller literature documents how human behavior adapts in response to AI quality. Dell'Acqua et al. (2023) study management consultants using AI assistance and find that AI-assisted workers progressively delegate cognitively demanding tasks to AI, producing measurable deskilling in unassisted task performance. Their effect sizes (0.3--0.5 SD) motivate the choice of |$\tau | \geq$ 0.2 as a minimum detectable effect threshold in the present study. Crucially, Dell'Acqua et al. measure task assignment behavior in a controlled experimental setting, not behavioral signals in naturalistic interaction logs.

Perez et al. (2023) demonstrate that AI sycophancy suppresses human expression of disagreement --- a finding complementary to BAA: if AI models are trained to suppress explicit corrections, measuring correction frequency as a BAA proxy will undercount genuine disengagement.

Shen et al. (2024) provide the most direct conceptual ancestor of the present work, surveying bidirectional human-AI alignment and identifying the measurement gap at the framework level. They propose behavioral proxy dimensions but provide no empirical measurement on public interaction logs.

\subsubsection{2.3 Temporal Analysis of AI Interaction Logs}

The LMSYS Chatbot Arena platform \citep{Zheng2023Judging} provides timestamped human preference votes across model pairs, making it the natural data source for temporal preference entropy analysis. The primary dataset requires approved research access; the publicly available fallback contains no timestamp field, making temporal binning impossible.

WildChat \citep{Zhao2024WildChat} provides approximately 1 million ChatGPT conversation logs with IP-hash pseudonymization, timestamps, and topic tags. Prior analyses of this dataset focus on content characterization, not on behavioral trend analysis within returning-user cohorts.

Mann-Kendall trend analysis (Mann, 1945; Kendall, 1975) is standard in environmental and behavioral time series analysis. The Hamed-Rao modification \citep{Hamed1998Modified} corrects for autocorrelated series; it is necessary in the present study given ACF lag-1 = 0.634 for the primary proxy series.

\subsubsection{2.4 Position}

Existing work on AI-to-human alignment does not account for human behavioral adaptation over time. Existing work on human behavioral adaptation uses controlled experiments, not public logs. Existing temporal analyses of WildChat and LMSYS Arena do not construct returning-user cohorts or test behavioral trends. Shen et al. (2024) identify the BAA framework conceptually but provide no empirical measurement. The present study occupies this gap.

